% The European Commission – Reading Summary
% Introduction to the European Union, Week 7
% Compile with: pdflatex european-commission-reading-summary.tex

\documentclass[11pt,a4paper]{article}
\usepackage[utf8]{inputenc}
\usepackage[T1]{fontenc}
\usepackage[margin=2.5cm]{geometry}
\usepackage{booktabs}
\usepackage{enumitem}
\setlist{nosep,leftmargin=*}
\usepackage{parskip}
\usepackage{microtype}
\usepackage{hyperref}
\usepackage{xcolor}
\definecolor{eublue}{RGB}{0,51,153}

\hyphenation{Spitzenkandidat Spitzenkandidaten}
\hypersetup{
  colorlinks=true,
  linkcolor=eublue,
  urlcolor=eublue,
  pdftitle={European Commission – Reading Summary},
  pdfauthor={Introduction to the EU, UCM}
}

\begin{document}

\title{The European Commission\\Key Themes \& Institutional Overview --- Rhetorical Analysis}
\author{Introduction to the European Union\\Universidad Complutense de Madrid\\Week 7 --- The Twin Executive (European Council \& Commission)}
\date{}
\maketitle
\vspace{-1em}
\hrule
\vspace{1em}

\section*{Rhetorical Framework (Ethos, Logos, Pathos)}

\begin{table}[h]
\centering
\begin{tabular}{@{}lll@{}}
\toprule
\textbf{Appeal} & \textbf{Meaning} & \textbf{Fact Sheet Use} \\
\midrule
\textbf{Ethos} & Credibility, authority & European Parliament Policy Department; treaty-based; official institutional source \\
\textbf{Logos} & Logic, evidence, argument & Treaty articles (TEU, TFEU); procedural detail; legal basis \\
\textbf{Pathos} & Emotion, stakes & Implicit: Commission as guardian of general interest; democratic legitimacy \\
\bottomrule
\end{tabular}
\end{table}

\textbf{Key insight:} The fact sheet presents the Commission in formal, legal-institutional terms. \textit{Ethos} comes from official EU documentation; \textit{logos} from systematic citation of treaty provisions; \textit{pathos} is understated but present in references to democratic scrutiny and citizens' interests.

\section*{Bibliographic Information}

\begin{itemize}
\item \textbf{Source:} European Parliament Fact Sheets on the European Union
\item \textbf{Fact Sheet:} ``The European Commission'' (Sheet 25)
\item \textbf{Authors:} Joanna Apap / Christophe Beaudouin (Policy Department for Justice, Civil Liberties and Institutional Affairs)
\item \textbf{Legal basis:} Art. 17 TEU; Arts. 234, 244--250, 290, 291 TFEU; Merger Treaty (1965)
\end{itemize}

\section*{Summary \& Key Points}

\subsection*{I. Core Role}

The Commission is the EU's \textbf{principal executive body} with:
\begin{itemize}
\item \textbf{Quasi-exclusive right of legislative initiative} --- proposes legislation
\item \textbf{Guardian of the Treaties} --- enforces EU law
\item \textbf{Manager} --- policies, budget, international agreements
\item \textbf{Regulator} --- competition, State aid, sanctions
\end{itemize}

\subsection*{II. History}

\begin{itemize}
\item \textbf{1951:} High Authority (ECSC) --- Jean Monnet model
\item \textbf{1957:} Separate Commissions for EEC and Euratom
\item \textbf{1965:} Merger Treaty --- single Commission of the European Communities
\item \textbf{2002:} ECSC Treaty expired; Commission inherited assets and research fund
\item \textbf{Lisbon (2009):} Commission president elected taking European election results into account
\end{itemize}

\subsection*{III. Composition (2024--2029)}

\begin{itemize}
\item \textbf{27 members:} 1 per Member State (Lisbon allows reduction to 2/3 of states, not yet applied)
\item \textbf{Structure:} President + 5 Executive Vice-Presidents + VP/HR (High Representative) + 20 Commissioners
\item \textbf{Appointment:} European Council proposes president (QMV) $\to$ Parliament elects $\to$ Council adopts list of Commissioners (with president-elect) $\to$ Parliament votes consent as body $\to$ European Council appoints
\item \textbf{Term:} 5 years, renewable (matches Parliament)
\end{itemize}

\subsection*{IV. Accountability}

\begin{itemize}
\item \textbf{Personal:} Commissioners must be independent; no other occupation; no instructions from governments (Art. 245 TFEU)
\item \textbf{Collective:} Motion of censure by Parliament $\to$ entire Commission resigns (Art. 234 TFEU). 16 motions tabled since 1972; none adopted
\item \textbf{Compulsory retirement:} Court of Justice may remove Commissioner for breach or serious misconduct (Art. 247 TFEU)
\end{itemize}

\section*{Main Arguments \& Powers}

\subsection*{I. Power of Initiative}

\textbf{A. Legislative}
\begin{itemize}
\item Proposes regulations and directives
\item Decision-making institutions (Parliament + Council) cannot adopt without a proposal; must base decision on proposal as presented
\end{itemize}

\textbf{B. Budgetary}
\begin{itemize}
\item Draws up draft budget (Art. 314 TFEU)
\item Own Resources Decision: Commission proposes; Council adopts unanimously after consulting Parliament
\item Receives estimates from all institutions and agencies; proposes contributions
\end{itemize}

\textbf{C. International}
\begin{itemize}
\item Negotiates agreements (Arts. 207, 218 TFEU) when Council gives mandate
\item Withdrawal (Art. 50): Commission submits recommendations on opening negotiations
\end{itemize}

\subsection*{II. Guardian of the Treaties}

\begin{itemize}
\item Ensures Treaties and secondary legislation are enforced
\item \textbf{Infringement procedure} (Art. 258 TFEU): Commission may bring Member State before Court of Justice for failure to fulfil obligations
\end{itemize}

\subsection*{III. Implementing Powers}

\begin{itemize}
\item \textbf{Competition:} Enforces rules; reviews State aid (Art. 108 TFEU)
\item \textbf{Budget:} Implements adopted budget; manages EU funds
\item \textbf{Crisis management:} Manages financial rescue packages; ESM voting procedure (with ECB)
\item \textbf{Delegated acts} (Art. 290 TFEU): Supplement/amend non-essential elements of legislation
\item \textbf{Implementing acts} (Art. 291 TFEU): Comitology replaced by advisory and examination procedures (Reg. 182/2011)
\end{itemize}

\subsection*{IV. Limited Initiative (Recommendations/Opinions)}

\begin{itemize}
\item \textbf{EMU:} European Semester assessments; exchange-rate recommendations; excessive deficit procedure; balance-of-payments assistance
\item \textbf{CFSP:} Supports VP/HR in submitting decisions to Council (Art. 30 TEU)
\end{itemize}

\section*{Organisation}

\begin{itemize}
\item \textbf{President:} Political guidance; internal organisation; allocates portfolios; appoints Vice-Presidents (with College approval)
\item \textbf{41 Directorates-General (DGs)} + Secretariat-General
\item \textbf{6 executive agencies} (e.g.\ Research, Climate, Health, Innovation, Education)
\item \textbf{Collegiality:} Acts by majority of members (Art. 250 TFEU); weekly meetings; written procedure for less sensitive matters; empowerment and sub-delegation for management
\end{itemize}

\section*{Relations with Parliament}

\begin{itemize}
\item \textbf{Principal interlocutor} in legislative and budgetary matters
\item \textbf{Discharge:} Parliament grants discharge to Commission (Art. 319 TFEU)
\item \textbf{Hearings:} Commissioners-designate subject to Parliament hearings; specific commitments tracked
\item \textbf{Right of initiative request} (Art. 225 TFEU): Parliament may request Commission to submit proposal
\item \textbf{Tensions:} Commission sometimes fails to comply with Parliament requests (e.g.\ Pegasus spyware) or delays reports; Parliament criticised Commission over EU-US data transfers (Schrems II, Data Privacy Framework)
\end{itemize}

\section*{Context \& Analysis}

\begin{itemize}
\item \textbf{Supranational core:} Commission embodies supranational logic --- acts in general interest, not national instructions
\item \textbf{Dual accountability:} To European Council (strategic) and Parliament (democratic)
\item \textbf{Expansion of roles:} From High Authority (coal/steel) to full policy range; agencies spun off; VP/HR double-hatted with Council
\item \textbf{Scrutiny gaps:} Parliament resolutions highlight Commission prioritising US relations over citizen data protection
\end{itemize}

\section*{Relevance to Course}

\begin{itemize}
\item \textbf{6th Session:} The Twin Executive --- Commission as one half of dual executive
\item \textbf{Institutional triangle:} Commission proposes; Parliament and Council decide
\item \textbf{Supranationalism:} Commission as engine of integration vs.\ intergovernmental Council/European Council
\item \textbf{Democratic legitimacy:} Hearings, censure, discharge, dialogue with Parliament
\end{itemize}

\section*{Discussion Questions}

\begin{enumerate}
\item Should the Commission be reduced to 18 members (2/3 of Member States) as Lisbon allows? What are the trade-offs?
\item How effective is Parliament's scrutiny of the Commission (hearings, censure, discharge)? Why have no motions of censure succeeded?
\item What is the significance of the Commission's ``quasi-exclusive'' right of initiative? How does it shape EU law-making?
\item How does the VP/HR's double role (Commission + Council/EEAS) affect the Commission's coherence in external relations?
\item Should the Commission have a dedicated Law Officer and Treasury Secretary, as reform proposals suggest?
\end{enumerate}

\vspace{1em}
\noindent\footnotesize\textit{Introduction to the European Union, Universidad Complutense de Madrid.}

\end{document}
